% Figure: the thermodynamic square (Born square / Guggenheim square), the
% mnemonic that encodes the four common potentials, their natural variables,
% and the Maxwell relations. Drawn with plain tikz nodes; no axis.
\begin{figure}[htbp]
\centering
\begin{tikzpicture}[font=\small]
  % corners of the square (variables) and edge midpoints (potentials)
  % layout: 4x4 unit square
  \def\s{3.6}
  % corner variables
  \node[engblack] (V)  at (0,\s)   {$V$};
  \node[engblack] (T)  at (\s,\s)  {$T$};
  \node[engblack] (S)  at (0,0)    {$S$};
  \node[engblack] (p)  at (\s,0)   {$p$};
  % edge-midpoint potentials
  \node[engblue]  (A)  at (0,{\s/2})    {$A$};
  \node[engblue]  (U)  at ({\s/2},\s)   {$U$};
  \node[engblue]  (G)  at (\s,{\s/2})   {$G$};
  \node[engblue]  (H)  at ({\s/2},0)    {$H$};
  % the square edges
  \draw[enggray, thick] (V) -- (T) -- (p) -- (S) -- (V);
  % the two diagonal arrows (sign convention of the square)
  \draw[engred, -{Stealth[length=2.2mm]}, thick] (S) -- (T);
  \draw[engred, -{Stealth[length=2.2mm]}, thick] (V) -- (p);
  % labels for the arrows
  \node[engred, font=\scriptsize, anchor=east] at ($(S)!0.5!(T)$) {};
\end{tikzpicture}
\caption[The thermodynamic square]{The thermodynamic square, a mnemonic for
the four common potentials and their natural variables. Each potential sits
between its two natural variables: internal energy $U$ between $S$ and $V$,
enthalpy $H$ between $S$ and $p$, Gibbs energy $G$ between $T$ and $p$, and
Helmholtz energy $A$ between $T$ and $V$. To read a differential, take the
potential and its two flanking corners; for example $U$ between $S$ and $V$
gives $\mathrm{d}U = T\,\mathrm{d}S - p\,\mathrm{d}V$, the conjugate
coefficient ($T$ for $S$, $p$ for $V$) read from the opposite corner along the
diagonal arrows, with the sign of the term set by the arrow direction (terms
pointing against an arrow carry a minus). The Maxwell relations are read off
the corners: equating the cross derivatives at any corner pair gives, for the
left and bottom edges, $(\partial S/\partial V)_T = (\partial p/\partial T)_V$.
The square stores the four potentials, their four differentials, and the four
Maxwell relations in one figure.}
\label{fig:vol04ch01:maxwell-square}
\end{figure}
